\hypertarget{security-engineering}{%
\chapter{SECURITY ENGINEERING}\label{security-engineering}}

\hypertarget{security-engineering-1}{%
\chapter{SECURITY ENGINEERING}\label{security-engineering-1}}

\hypertarget{fuzzing-libfuzzer-afl}{%
\subsection{36.1 Fuzzing (libFuzzer /
AFL++)}\label{fuzzing-libfuzzer-afl}}

Unit tests test what you \emph{know}. Fuzzing tests what you
\emph{don't}. * \textbf{Coverage-Guided:} The fuzzer instruments
binaries to see which inputs explore new code paths. *
\textbf{Sanitizers:} Always fuzz with AddressSanitizer (ASan) and
UndefinedBehaviorSanitizer (UBSan) enabled.

\hypertarget{cryptography-timing-attacks}{%
\subsection{36.2 Cryptography \& Timing
Attacks}\label{cryptography-timing-attacks}}

NEVER write your own crypto. Use \passthrough{\lstinline!libsodium!} or
\passthrough{\lstinline!BoringSSL!}. * \textbf{The Trap:}
\passthrough{\lstinline!memcmp(hash1, hash2, 32)!} exits early if the
first byte differs. * \textbf{The Attack:} Attacker measures time. If it
takes longer, they guessed the first byte right. * \textbf{The Fix:}
Constant-time comparison.
\passthrough{\lstinline"cpp     // libsodium's constant time check     if (crypto\_verify\_32(hash1, hash2) != 0) \{ /* Error */ \}"}

\hypertarget{side-channel-mitigations-spectre}{%
\subsection{36.3 Side-Channel Mitigations
(Spectre)}\label{side-channel-mitigations-spectre}}

Modern CPUs execute instructions speculatively. * \textbf{Scenario:}
\passthrough{\lstinline!if (x < array\_len) \{ val = array[x]; \}!} *
\textbf{Attack:} CPU predicts \passthrough{\lstinline!true!}, loads
\passthrough{\lstinline!array[x]!} (where \passthrough{\lstinline!x!} is
out of bounds secret). Even if checks fail,
\passthrough{\lstinline!array[x]!} is now in L1 cache. *
\textbf{Mitigation:} \passthrough{\lstinline!LFENCE!} (Load Fence) or
\passthrough{\lstinline!std::clamp!} indices to 0 on failure (masking).

\begin{center}\rule{0.5\linewidth}{0.5pt}\end{center}

\begin{center}\rule{0.5\linewidth}{0.5pt}\end{center}
